\documentclass[11pt]{article}
\usepackage[margin=2.5cm]{geometry}
\usepackage{amsmath, amssymb, amsthm}
\usepackage{bm}
\usepackage{hyperref}
\usepackage{array}

\newcommand{\R}{\mathbb{R}}
\newcommand{\dx}{\Delta x}
\newcommand{\dq}{\Delta q}

\title{P-RFO Notes\\
       \large Partitioned Rational-Function Optimization for Saddle Search}
\author{Notes for ChemDM}
\date{}

\begin{document}
\maketitle

\noindent A conceptual reference for partitioned rational-function optimization
as implemented in \texttt{src/chemdm/prfo.py}.  Captures the \emph{why}, not
just the \emph{what} --- in particular why BFGS fails for saddle search, where
the ``shift'' in P-RFO comes from, and why this is a saddle \emph{finder} and
not a reaction \emph{path} method.

\section{The problem}

Find a first-order saddle point of a potential energy surface $E(x)$ using
\emph{only forces} $-\nabla E$, no analytic Hessian.  The first-order saddle
has
\begin{equation*}
  \nabla E = 0, \qquad H = \nabla^2 E \text{ has exactly one negative eigenvalue}.
\end{equation*}

Two challenges: (a) the gradient alone is zero at every minimum \emph{and}
every saddle, so gradient-following cannot distinguish them, and (b) the
Hessian is expensive --- the very thing we want to avoid computing.

\section{Why BFGS specifically fails}

The BFGS quasi-Newton update
\begin{equation*}
  H_{k+1} = H_k + \frac{y y^T}{y^T s}
                - \frac{H_k s s^T H_k}{s^T H_k s},
  \qquad s = x_{k+1} - x_k, \quad y = g_{k+1} - g_k,
\end{equation*}
\emph{provably preserves positive definiteness} when the curvature condition
$y^T s > 0$ holds.  Line search enforces this condition for minimization.

The consequence for saddle search is fatal: if $H_0 = I$ (PD), every
subsequent $H_k$ is PD by induction.  The lowest eigenvalue can never cross
zero.  The Newton step $-H_k^{-1} g$ always points downhill.  \textbf{You
converge to a minimum, no matter how close you start to a saddle.}

This is not a corner case --- it is a fundamental structural mismatch.  BFGS
was designed for minimization and bakes that intent into its update rule.

The fix is to use a quasi-Newton update that \emph{does not preserve PD}:

\begin{center}
\begin{tabular}{l c l}
\hline
Update & PD-preserving & Used for \\
\hline
BFGS & yes & minimization \\
SR1 (symmetric rank-1) & no & saddle search (unstable if $r^T s$ small) \\
PSB (Powell-symmetric-Broyden) & no & saddle search (stable, slower) \\
\textbf{Bofill} (SR1/PSB convex mix) & no & saddle search (best of both) \\
\hline
\end{tabular}
\end{center}

\noindent \texttt{prfo.py} uses Bofill with weight
$\varphi = (r^T s)^2 / (\|r\|^2 \|s\|^2)$ --- the cosine squared between the
residual $r = y - H_k s$ and $s$.  SR1-heavy when those vectors align,
PSB-heavy when they do not.

\section{Newton with an indefinite Hessian --- almost works}

Suppose we \emph{had} the true indefinite Hessian at a saddle.  Decompose the
Newton step $\dx = -H^{-1} g$ in the eigenbasis:
\begin{equation*}
  \dx_i = -\frac{g_i}{\lambda_i}.
\end{equation*}

For positive $\lambda_i$, $\dx_i$ opposes $g_i$ --- descent.  For
\emph{negative} $\lambda_i$, $\dx_i$ aligns with $g_i$ --- \textbf{ascent}.
So Newton with the true Hessian automatically climbs the unstable mode and
descends the others.  This is eigenvector-following before we even introduce
RFO.

But two problems remain:
\begin{enumerate}
  \item \textbf{Step blows up near singular Hessians.}  Any $\lambda_i \approx 0$
        makes $|\dx_i| = |g_i / \lambda_i|$ enormous.  We need regularization.
  \item \textbf{Wrong-mode following.}  Early in the search, the ``naturally''
        most negative eigenvalue may not be the one we actually want to ascend
        (multiple soft modes, eigenvalues crossing zero during the search).
        We need explicit \emph{control} over which mode to follow.
\end{enumerate}

RFO solves both at once.

\section{The RFO step as shift-invert}

Replace the quadratic energy model with a \textbf{rational function}:
\begin{equation*}
  \Delta E(\dx) \approx
    \frac{g^T \dx + \tfrac{1}{2}\, \dx^T H \dx}{1 + \dx^T \dx}.
\end{equation*}

Stationary points of this rational function are solutions of the
\textbf{augmented eigenvalue problem}:
\begin{equation}\label{eq:augmented}
  \begin{pmatrix} H & g \\ g^T & 0 \end{pmatrix}
  \begin{pmatrix} \dx \\ 1 \end{pmatrix}
  = \mu
  \begin{pmatrix} \dx \\ 1 \end{pmatrix}.
\end{equation}

Expanding the first block row gives $H \dx + g = \mu \dx$, so
\begin{equation}\label{eq:shift-invert}
  \boxed{\;\dx = -(H - \mu I)^{-1} g.\;}
\end{equation}

\textbf{That is shift-invert.}  The Hessian gets shifted by $\mu$, inverted,
applied to $-g$.  The same operator appears in:
\begin{itemize}
  \item \textbf{Levenberg-Marquardt} for nonlinear least squares (shift =
        damping $\alpha$);
  \item \textbf{Trust-region Newton (Mor\'e-Sorensen)} --- shift chosen to
        satisfy $\|\dx\| = \text{trust\_radius}$;
  \item \textbf{Tikhonov regularization} for ill-posed inverse problems;
  \item \textbf{Inverse power iteration / shift-invert Lanczos} for eigenvalue
        problems.
\end{itemize}

What makes RFO distinctive among shift-invert methods is the \textbf{rule for
picking $\mu$}: it is an eigenvalue of the augmented $(N+1) \times (N+1)$
matrix in equation~\eqref{eq:augmented}.  Equivalently, substituting $\dx$
into the second row of \eqref{eq:augmented} gives the \textbf{secular
equation}:
\begin{equation}\label{eq:secular}
  \sum_i \frac{f_i^2}{\mu - \lambda_i} \;=\; \mu,
\end{equation}
where $\lambda_i$ are eigenvalues of $H$ and
$f_i = (\text{eigvec}_i)^T g$.

This scalar equation has $N+1$ roots, interleaved with the $\lambda_i$'s plus
one below the lowest and one above the highest eigenvalue.  \textbf{Each root
is a valid shift; each gives a different step type:}

\begin{itemize}
  \item \textbf{Lowest root} $\mu < \lambda_{\min}$: all $(\lambda_i - \mu) > 0$,
        so the step descends in every mode.  Minimization RFO.
  \item \textbf{Largest root} $\mu > \lambda_{\max}$: all $(\lambda_i - \mu) < 0$,
        so the step ascends in every mode.  Not usually what you want.
  \item \textbf{Intermediate roots}: mixed --- ascends in some modes, descends
        in others.
\end{itemize}

The choice of root encodes the method's intent.

The rational denominator $1 + \|\dx\|^2$ is what gives RFO its \textbf{automatic
step bounding}.  As $\|g\|$ grows or $H$ becomes singular, $\mu$ grows in
magnitude so $(\lambda_i - \mu)$ stays bounded away from zero.  The step $\dx$
cannot blow up by construction --- no trust-region constraint required.  (We
still impose one as a second safety net in \texttt{prfo.py}, used adaptively
via the $\rho = \Delta E_{\text{actual}} / \Delta E_{\text{predicted}}$ ratio.)

\section{The partitioning (the ``P'' in P-RFO)}

For transition-state search we want one mode to ascend and the rest to descend.
\textbf{Use two different shifts.}

\paragraph{Followed (ascent) mode $k$.}  Solve the $2 \times 2$ RFO subproblem
in mode $k$ alone:
\begin{equation*}
  \begin{pmatrix} \lambda_k & f_k \\ f_k & 0 \end{pmatrix}
  \begin{pmatrix} v \\ 1 \end{pmatrix}
  = \mu_+ \begin{pmatrix} v \\ 1 \end{pmatrix}.
\end{equation*}
Take the \emph{larger} root,
\begin{equation*}
  \mu_+ = \tfrac{1}{2}\!\left(\lambda_k + \sqrt{\lambda_k^2 + 4 f_k^2}\right)
        > \lambda_k.
\end{equation*}
Then $\dq_k = -f_k / (\lambda_k - \mu_+)$ ascends along mode $k$.

\paragraph{Descent subspace} (the other $N-1$ modes).  Solve the
$M \times M$ RFO subproblem on the remaining eigenvalues; take the \emph{lowest}
root $\mu_- < \min_{j \neq k} \lambda_j$.  Then
$\dq_j = -f_j / (\lambda_j - \mu_-)$ descends in each remaining mode.

Reassemble $\dq$ in the eigenbasis and rotate back to get $\dx$.  \textbf{One
ascent mode, all-others descent.}  That is the P in P-RFO.

\paragraph{Mode-following heuristic.}  The choice of which mode $k$ to follow:
on iteration 1, pick the lowest eigenvalue (most negative if you have one,
otherwise softest).  Subsequent iterations: pick the eigenvector with maximum
overlap with the previously-followed eigenvector.  Failure mode: when the
followed eigenvalue crosses zero or near-degenerate eigenvalues swap, overlap
can pick the wrong mode and the optimizer briefly walks the wrong way.  Our
HCN demo's iterations~19--31 show this in action --- the optimizer recovers,
but it is the main source of fragility in the current \texttt{prfo.py}.

\section{Coordinate systems are orthogonal to RFO}

RFO is just linear algebra on $(H, g)$.  It works in any coordinate system.
What changes between coordinate systems is the \emph{plumbing around} RFO, not
RFO itself.

\begin{center}
\begin{tabular}{l c c l}
\hline
Coords & Dim & Zero modes & Back-transform \\
\hline
Cartesian + trans/rot projection & $3N$ & $6$ (project out) & none \\
Non-redundant internals & $3N-6$ & $0$ & iterative Wilson B \\
Redundant internals (Pulay) & $M > 3N-6$ & $M - (3N-6)$ & Wilson B + pseudoinverse \\
\hline
\end{tabular}
\end{center}

The \textbf{three-axis design space} for any TS optimizer is:
\begin{quote}
\begin{tabular}{l l}
Step generator    & $\{$Newton, trust-region Newton, RFO/P-RFO, dimer rotation$\}$ \\
Hessian model     & $\{$analytic, Bofill, SR1, PSB, BFGS (minimization only)$\}$ \\
Coordinate system & $\{$Cartesian + projection, internal, redundant internal$\}$ \\
\end{tabular}
\end{quote}

Most published ``methods'' are just specific points in this cube.  Our
\texttt{prfo.py} is $\{$P-RFO, Bofill, Cartesian + projection$\}$.  Gaussian's
TS optimizer is $\{$P-RFO, Bofill, redundant internals$\}$.  Sella is roughly
$\{$trust-region Newton, exact Hessian via finite differences, Cartesian +
projection$\}$.  Dimer codes are $\{$dimer rotation step, no stored Hessian,
Cartesian$\}$.

These choices are independent --- you can swap any axis without changing the
others.

\section{P-RFO finds \emph{points}, not \emph{paths}}

The optimizer's iterates trace a discrete trajectory through configuration
space, but \textbf{that trajectory has no chemical meaning.}  It depends on
initial guess, Hessian model, trust radius, mode-following heuristic --- none
of which are properties of the molecule.  Two runs from different starts can
take very different optimizer paths to the same saddle.

The ``eigenvector-following'' name refers to following an \emph{eigenvector at
each iteration}, not to following a \emph{curve through configuration space}.
These are different objects.

The reaction path is a separate computation, defined intrinsically by the
PES:

\begin{description}
  \item[IRC (Intrinsic Reaction Coordinate, Fukui)] From the converged saddle,
    integrate $\mathrm{d}x/\mathrm{d}t = -g(x)$ in mass-weighted Cartesians,
    $\pm$ along the unstable mode.  One direction reaches the reactant
    minimum, the other the product minimum.  This curve is geometry-defined;
    different optimizers that find the same TS yield the same IRC.
\end{description}

This cleanly separates concerns:

\begin{center}
\begin{tabular}{l l l l}
\hline
Step & Method & Object computed & Cost \\
\hline
1 & P-RFO + Bofill & Saddle point (single geometry) & $\sim$20--100 force calls \\
2 & IRC (steepest descent $\pm$) & Reaction path (a curve) & $\sim$50--200 force calls \\
\hline
\end{tabular}
\end{center}

Both are needed for a complete reaction characterization.  P-RFO alone gives
you a saddle that might connect anywhere; IRC without a converged TS has no
starting point.

\section{A note on gradient extremals (and why we do not build them)}

\textbf{Gradient extremal curves} (Hoffman, Nord, Ruedenberg) are an alternative
class of methods that \emph{do} try to trace a curve directly from the reactant
basin.  A gradient extremal is the locus of points where $\nabla E$ is
parallel to a specific Hessian eigenvector.

There is a connection: the discrete trajectory of P-RFO with mode-following is
approximately a gradient extremal of the followed mode.  So in a loose sense,
P-RFO is discretely walking \emph{along} a gradient extremal toward its
endpoint at the saddle.

Gradient extremal \emph{following} methods (continuous ODE integration of the
curve, adaptive parametrization, branch tracking) make the \emph{curve itself}
the object of interest --- you walk from the reactant minimum along the curve,
through any saddles you encounter, and into the product basin.  This is
elegant and chemically motivated, but practically harder than the ``find TS
then IRC'' approach for several reasons:

\begin{itemize}
  \item Gradient extremals branch and recombine --- choosing the right branch
        at every bifurcation requires bookkeeping that IRC does not.
  \item They can leave the chemically relevant region of the PES, especially
        in high-dimensional systems.
  \item They do not necessarily pass through every saddle, and a saddle they
        do not pass through is invisible to the method.
  \item Numerically: ODE integration with branch tracking is heavier than
        ``find a stationary point, then steepest-descent.''
\end{itemize}

\textbf{Our practical choice is the simpler decomposition:}
\begin{center}
relaxed reactant
$\;\to\;$ \texttt{estimate\_lowest\_mode} (Lindh + Lanczos)
$\;\to\;$ perturb + P-RFO
$\;\to\;$ IRC.
\end{center}

The IRC step is cheap enough (gradient descent, $\sim$150 lines of code) that
we do not gain anything from trying to construct the reaction path directly
during the saddle search.  The gradient-extremal connection is a real
mathematical fact, but it does not pay off in code.

\section{Summary}

P-RFO is \textbf{shifted Newton with self-determined shifts}, where the shifts
come from an augmented eigenvalue problem and partition configuration space
into one ascent mode and a descent subspace.  The Hessian is approximated by
Bofill (an indefinite-capable quasi-Newton update --- BFGS does not work here).
The step is implicitly bounded by the rational-function model's denominator;
we add an explicit trust radius as a second safety net with adaptive
shrink/grow.

This finds a single saddle point per run.  To trace the reaction path that
connects the saddle to its endpoints, run IRC from the converged TS --- a
separate, cheap, well-defined computation.

For posterity: the three-axis design space (step generator $\times$ Hessian
model $\times$ coordinate system) is what makes the TS-optimization literature
navigable.  Once you see it that way, ``method X vs.\ method Y'' is mostly a
question of which point in the cube each method occupies.

\section{Postscript: minimization vs.\ TS is one root choice}

The cleanest possible statement of what RFO does: \textbf{minimization and
transition-state search differ only by which root of the same secular equation
you pick.}  Everything else --- Hessian update, Bofill, trust radius,
coordinate system --- is shared infrastructure.

Recall the secular equation~\eqref{eq:secular}:
\begin{equation*}
  \sum_i \frac{f_i^2}{\mu - \lambda_i} = \mu.
\end{equation*}
By the interlacing property of its $N+1$ roots, they sit between the Hessian
eigenvalues plus one outside on each end:
\begin{equation*}
  \mu_0 \;<\; \lambda_1 \;<\; \mu_1 \;<\; \lambda_2 \;<\; \mu_2 \;<\; \lambda_3
  \;<\; \cdots \;<\; \lambda_N \;<\; \mu_N.
\end{equation*}

Each choice of root gives a step of the form $\dx = -(H - \mu I)^{-1} g$.
What changes is the \emph{sign pattern} of $(\lambda_i - \mu)$ across modes:

\begin{center}
\begin{tabular}{l l l}
\hline
Root chosen & Sign pattern of $(\lambda_i - \mu)$ & Step type \\
\hline
$\mu_0 < \lambda_1$            & all positive             & minimization \\
$\mu_1 \in (\lambda_1, \lambda_2)$ & one negative, rest positive  & TS, follow mode 1 \\
$\mu_2 \in (\lambda_2, \lambda_3)$ & two negative, rest positive  & 2nd-order saddle search \\
\hline
\end{tabular}
\end{center}

So, to answer the practical question: \textbf{starting from a local minimum,
pick $\mu = \mu_1$ instead of $\mu_0$ and the optimizer automatically pursues
a transition-state search along the softest mode.}  No fake-negative-eigenvalue
initialization required, no special-case code, no different update rule.
Same Hessian, same gradient, same secular equation --- different root.

Two practical caveats:

\begin{enumerate}
  \item At a true local minimum, $\|g\| = 0$ exactly, so all $f_i = 0$ and
        the secular equation degenerates: every $\mu_i = \lambda_i$, every step
        is zero.  You must perturb slightly off the minimum to give the
        optimizer something to climb (and to break mode degeneracy in
        symmetric systems).  This is what the demo's
        \texttt{x\_start = x\_min + kick * u} step does.
  \item The two-shift formulation in Section~5 ($\mu_+$ from a $2 \times 2$
        subproblem on the followed mode, $\mu_-$ from an $M \times M$
        subproblem on the rest) is mathematically equivalent to picking
        $\mu_1$ here.  The partitioned form is what production codes
        implement for numerical reasons (smaller eigenproblems, cleaner
        mode-following control), but the underlying mathematical content is
        ``pick the secular root that sits above the eigenvalue you want to
        ascend.''
\end{enumerate}

This is the unifying picture: P-RFO and ordinary RFO minimization are the
\emph{same algorithm} running on the \emph{same data structure}, choosing a
different root of the same scalar equation.  When you see it that way, the
whole machinery of TS search collapses into a one-line modification of an
ordinary minimizer.

\section*{References}

\begin{description}
\item[Schlegel,] \emph{WIREs Comput.\ Mol.\ Sci.} \textbf{1}, 790 (2011).
      Best tutorial review.
\item[Banerjee, Adams, Simons, Shepard,] \emph{J.\ Phys.\ Chem.} \textbf{89},
      52 (1985).  The original P-RFO paper.
\item[Baker,] \emph{J.\ Comput.\ Chem.} \textbf{7}, 385 (1986).  The
      practical recipe (mode following by max overlap) most QC codes follow.
\item[Bofill,] \emph{J.\ Comput.\ Chem.} \textbf{15}, 1 (1994).  The Bofill
      update.
\item[Besal\'u \& Bofill,] \emph{Theor.\ Chem.\ Acc.} \textbf{100}, 265
      (1998).  Proper trust-radius treatment for RFO.
\item[Wales,] \emph{Energy Landscapes} (Cambridge, 2003), ch.\ 6--7.
      Wide-angle view connecting RFO to dimer, GAD, and gradient extremals.
\item[Hoffman, Nord, Ruedenberg,] \emph{Theor.\ Chim.\ Acta} \textbf{69}, 265
      (1986).  Gradient extremal curves (for reference; we do not implement
      these).
\end{description}

\end{document}
